\documentclass{article}
\usepackage[T1]{fontenc}
\usepackage{lmodern}
\usepackage{graphicx}
\usepackage{amsmath,amssymb}
\usepackage[a4paper,margin=2cm]{geometry}
\usepackage[absolute,overlay]{textpos}
\setlength{\TPHorizModule}{1mm}
\setlength{\TPVertModule}{1mm}
\usepackage[indonesian]{babel}
\usepackage{hyperref}
\setlength{\emergencystretch}{2em}
% unit: Halaman 22

\begin{document}

% === Halaman 22 (llm) ===

\begin{itemize}
    \item \texttt{XLEX} dipetakan menjadi \texttt{th*t}.
    \item Kata yang cocok untuk \texttt{th*t} adalah \texttt{that}.
    \item Jadi kita memperoleh: $\text{E} \rightarrow \text{a}$
    \item Hasil iterasi pertama:
\end{itemize}

\vspace{5mm}

\begin{center}
\texttt{heVeTCSWPeYVaWHaVSReQMthaYVaOeaWHRtatePFaMVaWHKVSTYhtZe\\
theKeetPeJVSZaYPaRRGaReMWQhMGHtQaReWGPSReHMtQaRaKeaTtM\\
JTPRGaVaKaeTRaWHattattMZeTWAWSQWtSwatTVaPMRtRSJGSTVRea\\
YVeatCVMUeMWarGMeWtMJMGCSMWtSJOmeQtheVeQeVetQSVSTWHKPaG\\
ARCStRWeaVSWeeBtVeZMtFSJtheKGaAAWhaPSWYSWeWeaVtheStheVt\\
heRGaPeRQeVeeBGeeHMWYPFhaVHAWHYPSR RFQMthaPPtheaCCeaVaWG\\
eSJKTVWMRheHYSPHtheQeMYhtSJtheMWReGtQaROeVFeZaVAaKPeaW\\
HtaAMWYaPPthMWYRMWtSGSWRMHeVatMSWMGSTPHhaVHPFKPaZeNTCMt\\
eVJSVhMRSCMWMSWVeRCeGtMWYyMt}
\end{center}

\end{document}
